% Lösungen Einheit 2 · Ausklammern
\zweigkopf{Einheit 2 von 5 · Ausklammern}{Hier lernst du, einen gemeinsamen Faktor vor die Klammer zu ziehen · neu oder schon bekannt – je nach Buch · keine P10-Aufgabe · baut auf: Zahl mal Klammer (Einheit 1)}{l2}
\erg{16}{a) $3 \cdot (x + 4)$ \quad b) $5 \cdot (m + 2)$ \quad c) $7 \cdot (y + 3)$ \quad d) $2 \cdot (x - 8)$ \quad e) $x \cdot (x + 9)$ \quad f) $4x \cdot (x + 3)$ \quad g) $5x \cdot (x + 1)$ \quad h) $3 \cdot (2x^2 + x - 3)$ \quad i) $5x \cdot (2x - 3)$; Probe: $10x^2 - 15x$ \quad j) $4m \cdot (3m + 2n - 1)$; Probe: $12m^2 + 8mn - 4m$}
\erg{17}{a) Ben hat den Faktor $6$ nur aus dem ersten Glied gezogen; richtig: $6 \cdot (x + 3)$ \quad b) $8 \cdot (m - 3)$ \quad c) Ella hat die $1$ in der Klammer vergessen ($7y : 7y = 1$); richtig: $7y \cdot (y + 1)$ \quad d) $9x \cdot (x - 1)$}
\erg{18}{a) ja, $2$ \quad b) nein \quad c) nein \quad d) ja, $4m$ \quad e) nicht falsch, aber nicht vollständig: in der Klammer steckt noch der Faktor $2$; ganz ausgeklammert $4 \cdot (3x + 2)$}
\erg{19}{a) $6k + 15 = 3 \cdot (2k + 5)$, Länge $(2k + 5)$~cm \quad b) $k^2 + 8k = k \cdot (k + 8)$, Länge $(k + 8)$~cm \quad c) Länge $13$~cm; $3 \cdot 13 = 39$ und $6 \cdot 4 + 15 = 39$}
